\documentclass[11pt,a4paper]{article}
\usepackage[margin=1in]{geometry}
\usepackage{hyperref}
\usepackage{enumitem}
\usepackage{graphicx}
\usepackage{xcolor}

% -----------------------------------------------------
% bvraghav-tiet-ucs503p-article.sty
% -----------------------------------------------------

% Variable: Lab Instructor
\makeatletter \newcommand{\BvrSetLabInstructor}[1]{%
  \gdef\Bvr@LabInstructor{#1}}
% default
\newcommand{\Bvr@LabInstructor}{Dr. Raghav %
  B. Venkataramaiyer}
% public accessor
\newcommand{\BvrLabInstructorValue}{\Bvr@LabInstructor}
\makeatother

% Add subtitle
\usepackage{titling}
% -- Set title bold
\pretitle{\begin{center}\bfseries\fontsize{18bp}{18bp}\selectfont}
\makeatletter
\newcommand{\BvrSubtitle}[1]{%
  \posttitle{%
    \par\end{center}
    \begin{center}\large#1\end{center}
    \vskip 0.5em}%
}
\makeatother

% Hints in the section title(s)
\newcommand{\BvrHint}[1]{%
  {\normalfont\normalsize\itshape\hfill #1}}

% -----------------------------------------------------

\title{Airport Management System}

\BvrSetLabInstructor{Dr. Raghav B. Venkataramaiyer}

\BvrSubtitle{%
  Project Proposal for UCS503P  \\
  Submitted to: \BvrLabInstructorValue \\ \bfseries
  Thapar Institute of Engineering and Technology%
}

\author{
  TODO Name, CSED%
  \thanks{Roll No: \texttt{10xxxxxxxx}, %
    Email: \texttt{TODO\@thapar.edu}}
}

\date{\today}

\begin{document}
\maketitle

\section{Higher-order goal%
  \BvrHint{\ldots secondary framing}}
This project aims to demonstrate that a single, well-designed
operational data layer -- shared across every subsystem of an airport
-- can support a wide range of day-to-day airport functions (flight
operations, passenger-facing services, staff management, security,
and facilities) without each subsystem owning a private data silo.
The immediate engineering objective is to translate this architectural
claim into a deployable, multi-module web platform.

\section{Time-to-value %
\BvrHint{\ldots primary heuristic}}
\begin{itemize}[leftmargin=0.7cm]
    \item We will deliver a live, working skeleton (database + API +
      frontend health check) before any feature work begins, so every
      subsequent module ships against a real, deployed system.
    \item Modules are explicitly tiered: \textbf{Tier 1} modules
      (concurrency, cascading updates, real-time state, prediction)
      are built first for engineering depth; \textbf{Tier 2} modules
      (CRUD + status pipelines) follow once the core patterns are
      proven, so complexity is front-loaded by design rather than
      discovered under deadline pressure.
    \item Success in each weekly iteration is defined by what is
      deployed and demonstrable that week, not by a big-bang release
      at the end.
\end{itemize}

\section{Problem Statement}
Airports run many operationally distinct but data-coupled functions
-- flight scheduling, gate assignment, baggage handling, passenger
booking, staff rostering, security screening, parking, and facilities
-- that in practice are often built as separate systems stitched
together after the fact. This causes:
\begin{itemize}[leftmargin=0.7cm]
    \item \textbf{Data duplication and drift:} the same flight, gate,
      or passenger record is re-entered or re-synced across systems,
      creating inconsistency.
    \item \textbf{Delayed cascades:} a single flight delay should
      ripple through gate reassignment, connecting-flight risk,
      staff scheduling, and passenger notifications -- but siloed
      systems propagate this slowly or not at all.
    \item \textbf{Concurrency hazards:} shift scheduling, resource
      assignment, and parking-slot booking are all instances of the
      same conflict-detection problem, usually re-solved (and
      re-bugged) independently in each silo.
    \item \textbf{Fragmented staff tooling:} ground crew, security,
      medical, and admin staff each need role-specific views into the
      same underlying operational state.
\end{itemize}

\textbf{Why this matters now (contextual relevance):} even a
teaching-scale airport platform surfaces the same architectural
pressure real systems face -- one shared operational core, read and
written by many role-specific modules, is what actually prevents the
inconsistency and cascade-delay problems above.

\section{Proposed Solution %
\BvrHint{\ldots Software Proposition}}
\subsection{Overview}
Build a \textbf{web-based} Airport Management System around one
shared operational core (\texttt{Airport}, \texttt{Gate},
\texttt{Flight}, \texttt{Passenger}, \texttt{Resource}, \texttt{User}),
with the following module groups on top of it:
\begin{itemize}[leftmargin=0.7cm]
    \item \textbf{Ops Core:} flights, gates, resources, turnaround,
      and delay handling with cascading recalculation.
    \item \textbf{Passenger-facing:} booking portal, check-in,
      connecting-flight risk, baggage tracking, accessibility \&
      medical assistance requests, a support chatbot, parking, and
      a facilities directory.
    \item \textbf{Operations \& safety:} emergency handling, security
      checkpoints/screening/incidents, predictive analytics for delay
      and congestion.
    \item \textbf{Staff:} Staff \& Employee Management -- records,
      shift scheduling with conflict prevention, and attendance.
    \item \textbf{Admin Dashboard:} rolls up all of the above.
\end{itemize}

\subsection{Core workflow}
\begin{enumerate}[leftmargin=0.7cm]
    \item Ops staff maintain flights, gates, and resources against the
      shared core; a delay on any flight triggers cascading updates
      to dependent gate/resource/connecting-flight state.
    \item Passengers interact with the same core through booking,
      check-in, baggage status, and the chatbot -- all reading live
      operational data rather than a stale copy.
    \item Staff are scheduled onto shifts through a single
      conflict-checking service; the same pattern is reused for
      resource assignment and parking-slot booking.
    \item The Admin Dashboard aggregates status and metrics across
      every module from the shared core.
\end{enumerate}

\subsection{Operational constraints for an educational setting}
\begin{itemize}[leftmargin=0.7cm]
    \item Role-based access (passenger, ground\_crew, security,
      medical, ops\_manager, admin) enforced uniformly via JWT
      middleware rather than per-module ad hoc checks.
    \item Avoid dependence on external/paid services where a mock or
      rule-based substitute suffices (e.g.\ mock payments, rule-based
      chatbot over live DB data) -- focus stays on workflow and system
      correctness.
    \item Design for deployability using common web hosting and a
      managed Postgres instance.
\end{itemize}

\section{Solution Approach %
\BvrHint{\ldots Engineering focus}}
This is an engineering course project, so the focus is on the shared
data layer, correctness of concurrent operations, and workflow --
not on research-grade algorithms.

\subsection{Conflict detection %
\BvrHint{\ldots reused pattern, not novel research}}
Shift scheduling, resource assignment, and parking-slot booking all
reduce to the same overlap-detection problem over time ranges:
\begin{itemize}[leftmargin=0.7cm]
    \item A single, unit-tested conflict-checking service
      (\texttt{shiftConflict.js}) is written once against shifts and
      reused as-is for resource assignment and parking.
    \item On conflict, the API returns \texttt{409} with the
      conflicting rows so the caller (staff UI) can resolve it, rather
      than silently overwriting or rejecting without explanation.
\end{itemize}

\subsection{Web architecture %
\BvrHint{\ldots web-based solution}}
A standard 3-tier web architecture:
\begin{itemize}[leftmargin=0.7cm]
    \item \textbf{Frontend:} React (Vite), role-aware views for
      passengers, staff, and admins.
    \item \textbf{Backend API:} Node.js + Express, JWT auth,
      Socket.IO for real-time updates (delays, gate changes).
    \item \textbf{Database:} PostgreSQL, one schema covering all
      15 modules' entities against the shared operational core.
\end{itemize}

\subsection{Testing and delivery}
\begin{itemize}[leftmargin=0.7cm]
    \item Jest (backend) for unit tests on pure logic (e.g.\ conflict
      detection) that don't require a live database; React Testing
      Library planned for frontend component tests.
    \item Each week's module ships against the already-deployed
      skeleton, so the live system is always in a demonstrable state.
\end{itemize}

\section{Evaluation Criterion %
\BvrHint{\ldots measurable and attributable}}

\subsection{Primary evaluation metric}
\textbf{Correctness of shared-core consistency:} whether a state
change in one module (e.g.\ a flight delay) is correctly and promptly
reflected in every dependent module (gate status, connecting-flight
risk, passenger notification) without manual re-entry.
\begin{itemize}[leftmargin=0.7cm]
    \item Target: cascading updates visible across dependent modules
      within the same request/socket cycle that recorded the delay.
    \item Attribution: verified via integration tests and manual
      cross-module checks during weekly demos.
\end{itemize}

\subsection{Secondary evaluation metrics}
\begin{itemize}[leftmargin=0.7cm]
    \item \textbf{Conflict-detection correctness:} unit-test coverage
      of the shared overlap-checking service across shifts, resource
      assignment, and parking.
    \item \textbf{API health:} \texttt{/health} reports
      \texttt{db.connected: true} in every deployed environment.
    \item \textbf{Role isolation:} each protected route enforces the
      correct role set, verified by auth middleware tests.
    \item \textbf{Weekly deployability:} each week's increment builds
      and deploys without breaking the previous week's live system.
\end{itemize}

\section{Scalability %
\BvrHint{\ldots with foresight}}
\begin{enumerate}[leftmargin=0.7cm]
    \item \textbf{Scalable (theoretically):} decoupled frontend/backend,
      stateless JWT-authenticated API, indexed queries on the shared
      core tables, and a conflict-checking service designed to be
      reused (not re-implemented) by every module that needs it.
    \item \textbf{Operationally deployable (practically):} managed
      Postgres, platform-hosted backend (Render/Railway) and frontend
      (Vercel), configuration entirely via environment variables.
\end{enumerate}

\section{Engine availability heuristic}
\begin{itemize}[leftmargin=0.7cm]
    \item Web framework (Express) for REST APIs and Socket.IO for
      real-time push.
    \item Managed relational database (Postgres via
      Render/Neon/Supabase).
    \item Token-based authentication (JWT) with role middleware.
    \item Test framework (Jest / React Testing Library).
    \item Platform-hosted CD targets (Render/Railway, Vercel) instead
      of self-managed infrastructure.
\end{itemize}
These mature components let the project focus on the shared data
model, cascading/conflict logic, and role-based workflow rather than
infrastructure.

\section{Project scope and deliverables %
\BvrHint{\ldots iteration-friendly}}
\subsection{Delivered so far %
\BvrHint{\ldots Weeks 1--2}}
\begin{itemize}[leftmargin=0.7cm]
    \item Week 1: finalized scope and tiering across 15 modules, full
      ERD, complete Postgres schema, backend/frontend skeleton with a
      live \texttt{/health} check.
    \item Week 2: JWT auth (register/login/me), Staff \& Employee
      Management (records, shift scheduling with conflict rejection,
      unit tested), attendance clock-in/out, role-aware frontend
      dashboard.
\end{itemize}

\subsection{Subsequent deliverables}
\begin{itemize}[leftmargin=0.7cm]
    \item Ops Core (flights, gates, resources, turnaround) and delay
      handling with cascading recalculation.
    \item Booking portal, connecting-flight risk, baggage management,
      check-in.
    \item Emergency handling, security module, predictive analytics,
      support chatbot.
    \item Parking, facilities directory, accessibility \& medical
      assistance.
    \item Admin dashboard rolling up all modules.
\end{itemize}

\section{Risks and mitigations}
\begin{itemize}[leftmargin=0.7cm]
    \item \textbf{Shared-core coupling risk:} a bug in the core schema
      affects every module; mitigated by finalizing the ERD and schema
      in Week 1, before feature work, and unit-testing shared services
      like conflict detection.
    \item \textbf{Scope size (15 modules):} mitigated by the explicit
      Tier 1/Tier 2 split, so depth is spent where it matters and the
      remainder is deliberately kept simple.
    \item \textbf{Real-time correctness:} cascading updates (delays,
      gate changes) are harder to get right than simple CRUD;
      mitigated by building this pattern early (Ops Core, Tier 1) and
      reusing the same conflict-service pattern established in Week 2.
    \item \textbf{Deployment drift:} mitigated by deploying the
      skeleton in Week 1 and shipping every subsequent module against
      that same live environment.
\end{itemize}

\section{Summary}
This project proposes a deployable, multi-module Airport Management
System built on one shared operational data layer, demonstrating that
a well-designed core prevents the data duplication and cascade-delay
problems that fragmented airport systems typically suffer from. It
meets the course criteria by shipping a deployed skeleton in Week 1,
tiering module complexity deliberately, reusing a single unit-tested
conflict-detection pattern across multiple modules, and defining
cross-module consistency as the primary, attributable evaluation
metric.

\end{document}
